\cvsection{\tr{Esperienze Lavorative}{Work Experience}}
\smallskip
\cvcompany{SCM Group S.p.A.}{Rimini}{\tr{Ott 2023 – Presente}{Oct 2023 – Present}}{\cvtag{Python}\cvtag{RAG}\cvtag{LLM Agents}\cvtag{Azure OpenAI}\cvtag{Vertex AI}\cvtag{MongoDB}\cvtag{Fine-Tuning}\cvtag{Prompt Engineering}\cvtag{LLM-as-a-Judge}}
\begin{cvroles}
\cvrole{AI Specialist}{\tr{Gen 2026 – Presente}{Jan 2026 – Present}}
\begin{cvitems}
  \item \tr{Progettata e sviluppata la parte \textbf{RAG} (ingestion e retrieval) di un sistema ad \textbf{agenti LLM} costruito da zero in team, per onboarding, hotline e ricambi: \textbf{oltre 500.000 documenti}, oggi in \textbf{beta interna} su 5 settori con \textbf{circa 100 utenti} in crescita.}
           {Designed and built the \textbf{RAG} layer (ingestion and retrieval) of an \textbf{LLM agent} system developed from scratch by the team, for onboarding, hotline and spare parts: \textbf{500,000+ documents}, now in \textbf{internal beta} across 5 areas with \textbf{about 100 users} and growing.}
  \begin{cvitems}
    \item \tr{\textbf{Ingestion}: ticket di assistenza e ricambi, manualistica; PDF legacy convertiti in formato strutturato con Google Document AI.}
             {\textbf{Ingestion}: support tickets, spare-parts tickets and manuals, with legacy PDFs converted into structured data with Google Document AI.}
    \item \tr{\textbf{Retrieval}: prefiltraggio in MongoDB tramite tool degli agenti, su campi descritti da \textbf{ontologie} di settore; ricerca vettoriale su Vertex AI e reranking.}
             {\textbf{Retrieval}: MongoDB pre-filtering via agent tools on fields defined by area \textbf{ontologies}; Vertex AI vector search and reranking.}
    \item \tr{\textbf{Sicurezza}: ogni ricerca riguarda solo i documenti a cui l'utente ha accesso, e il perimetro lo decide il codice, non il modello, così i settori non si mescolano.}
             {\textbf{Security}: searches cover only documents the user can access; the scope is set by code, not by the model, so areas never mix.}
    \item \tr{\textbf{Rollout per settore}: interviste sui pain point, valutazione di fattibilità, ontologie e agenti dedicati, test con \textbf{LLM-as-a-judge} e sviluppo guidato dal feedback.}
             {\textbf{Rollout by business area}: pain-point interviews, feasibility assessment, dedicated ontologies and agents, testing with \textbf{LLM-as-a-judge} and feedback-driven development.}
  \end{cvitems}
  \item \tr{Sviluppato un servizio di \textbf{traduzione di dominio} basato su LLM, con fine-tuning e prompt engineering, iniziato nel ruolo precedente.}
           {Built a \textbf{domain translation} service on LLMs, with fine-tuning and prompt engineering, started in the previous role.}
\end{cvitems}
\smallskip
\cvrole{Data Analyst/Developer}{\tr{Apr 2024 – Dic 2025}{Apr 2024 – Dec 2025}}
\begin{cvitems}
  \item \tr{Sviluppati algoritmi di \textbf{manutenzione predittiva} su dati in tempo reale di macchinari industriali.}
           {Developed \textbf{predictive maintenance} algorithms on real-time data from industrial machinery.}
  \item \tr{Progettata una soluzione rivolta ai clienti, oggetto di una \textbf{domanda di brevetto} in corso.}
           {Designed a customer-facing solution that is the subject of a pending \textbf{patent application}.}
\end{cvitems}
\smallskip
\cvrole{\tr{Data Analyst – Tirocinio Curriculare}{Data Analyst Intern}}{\tr{Ott 2023 – Feb 2024}{Oct 2023 – Feb 2024}}
\begin{cvitems}
  \item \tr{Sviluppato un sistema di \textbf{classificazione non supervisionata} di documenti tecnici con topic modeling (LDA e BERTopic), oggetto della tesi di laurea.}
           {Built an \textbf{unsupervised classification} system for technical documents with topic modeling (LDA and BERTopic), the subject of my bachelor's thesis.}
  \item \tr{Addestrati modelli di \textbf{Machine Learning supervisionato} sulla piattaforma \href{https://www.expert.ai/}{\textit{Expert.ai}}.}
           {Trained \textbf{supervised Machine Learning} models on the \href{https://www.expert.ai/}{\textit{Expert.ai}} platform.}
\end{cvitems}
\end{cvroles}

\medskip
\cvrow{\small\color{emphasis}\textbf{\tr{Frontend Developer – Tirocinio Curriculare}{Frontend Developer Intern}}, \tr{Università di Bologna}{University of Bologna}, Cesena}{\small\faCalendar~\tr{Set 2020 – Nov 2020}{Sep 2020 – Nov 2020}}
\smallskip
\cvrow{\small\color{emphasis}\textbf{Barista}, Baruffa Caffè, Milano Marittima (\tr{stagionale, durante gli studi}{seasonal, while studying})}{\small\faCalendar~Mar 2015 – Nov 2020}
